\begin{definition}[Haar-integral local constant]\label{mod:delta-integraldefinition}
Let $F$ be a nonarchimedean local field.  Let
$\chi:F^\times\to\mathbf C^\times$ have exact multiplicative conductor
$m\in\mathbf N$, and let $\psi:F\to\mathbf C^\times$ have exact additive
conductor $n\in\mathbf Z$ in the conventions of
Definition~\ref{mod:basic-characterconductors}: $m$ is the least depth at
which $\chi$ is trivial on $U_F^m$, while $\mathfrak p_F^{-n}$ is the
largest fractional ideal on which $\psi$ is trivial.  An
\emph{integral-admissible denominator} is a proof-bearing element
$\gamma\in F^\times$ satisfying
\[
  \operatorname{ord}_F(\gamma)=m+n.
\]
Equivalently, $\gamma\mathcal O_F=\mathfrak p_F^{m+n}$.  Such
denominators exist.  If $\gamma$ and $\gamma'$ are admissible, the
explicit ratio $a=\gamma'/\gamma$ belongs to
$U_F^0=\mathcal O_F^\times$ and satisfies $\gamma'=a\gamma$.

Lean equips the subtype $U_F^0$ with its Borel measurable structure and
proves that it is a compact multiplicative group.  Let $du$ be an arbitrary
positive left Haar measure on this group: formally it is a measure with
Mathlib's \texttt{IsHaarMeasure} instance, which includes left invariance,
positive mass on every nonempty open set, and finite mass on compact sets.
In particular,
\[
  0<du(U_F^0)<\infty.
\]
No probability normalization or prescribed value of $du(U_F^0)$ is chosen.

For $u\in U_F^0$, define exactly
\[
  g_\gamma(u)=\psi(u/\gamma)\chi(u)^{-1},
  \qquad
  I_{du,\gamma}(\chi,\psi)
    =\int_{U_F^0}g_\gamma(u)\,du.
\]
The Lean integral is over the whole subtype $U_F^0$ itself, not over the
additive field, all of $F^\times$, or a set with an additive Haar measure.
The function $g_\gamma$ is continuous and Borel measurable, and it is
integrable for every such Haar measure.  Rescaling a measure by
$c\in\mathbf R_{\geq0}$ gives the literal identity
\[
  I_{c\,du,\gamma}=c\,I_{du,\gamma}.
\]

Using the total phase of Definition~\ref{mod:basic-phase}, with
$\ph(0)=1$ and $\ph(z)=z/|z|$ for $z\ne0$, define
\[
  \Delta^{\mathrm{int}}_{F,du}(\chi,\psi;\gamma)
    =\chi(\gamma)\,\ph\!\left(I_{du,\gamma}(\chi,\psi)\right).
\]
For any two positive left Haar measures $du$ and $dv$, Haar uniqueness on
the compact group writes $dv=c\,du$ for a strictly positive real scalar.
Positive real scaling does not change the total phase, so Lean proves
$\Delta^{\mathrm{int}}_{F,dv}=\Delta^{\mathrm{int}}_{F,du}$ without
renormalizing either measure.

The admissible-denominator transformation has the manuscript's exact
orientation:
\[
  g_{\gamma'}(au)=\chi(a)^{-1}g_\gamma(u),
  \qquad
  I_{du,\gamma'}=\chi(a)^{-1}I_{du,\gamma}.
\]
The second equality uses left Haar invariance under the substitution
$u=av$.  Exact multiplicative-conductor data imply
$|\chi(a)|=1$; the formal proof factors the restriction of $\chi$ through
the finite conductor quotient only for this unitary assertion and never
chooses a quotient representative.  Consequently, if
$I_{du,\gamma}\ne0$, Lean proves that the two admissible denominators give
the same integral local constant.  The nonzero hypothesis is stated
explicitly because the total phase has a separate value at zero; this node
does not assume the later integral--finite bridge or its nonvanishing
consequence.

\LeanFile{LanglandsFirstMainLemma/Delta/IntegralDefinition.lean}
\lean{LanglandsFirstMainLemma.IntegralAdmissibleGamma}
\lean{LanglandsFirstMainLemma.IntegralAdmissibleGamma.exists_admissible}
\lean{LanglandsFirstMainLemma.IntegralAdmissibleGamma.smul_lattice_zero}
\lean{LanglandsFirstMainLemma.IntegralAdmissibleGamma.ratio_mem_unitFiltration}
\lean{LanglandsFirstMainLemma.isOpen_unitFiltration}
\lean{LanglandsFirstMainLemma.isOpen_unitFiltrationInside_zero}
\lean{LanglandsFirstMainLemma.unitFiltrationZeroMeasurableSpace}
\lean{LanglandsFirstMainLemma.unitFiltrationZeroCompactSpace}
\lean{LanglandsFirstMainLemma.integralQuasiChar_norm_eq_one_of_mem_unitFiltration_zero}
\lean{LanglandsFirstMainLemma.unitHaarIntegrand}
\lean{LanglandsFirstMainLemma.unitHaarIntegrand_continuous}
\lean{LanglandsFirstMainLemma.unitHaarIntegrand_measurable}
\lean{LanglandsFirstMainLemma.unitHaarIntegrand_integrable}
\lean{LanglandsFirstMainLemma.unitHaarMeasure_univ_pos}
\lean{LanglandsFirstMainLemma.unitHaarMeasure_univ_lt_top}
\lean{LanglandsFirstMainLemma.unitGaussIntegral}
\lean{LanglandsFirstMainLemma.unitGaussIntegral_smul_nnreal_measure}
\lean{LanglandsFirstMainLemma.unitGaussIntegral_admissibleGamma_change}
\lean{LanglandsFirstMainLemma.deltaIntegral}
\lean{LanglandsFirstMainLemma.deltaIntegral_measure_independent}
\lean{LanglandsFirstMainLemma.deltaIntegral_gamma_independent_of_ne_zero}
\leanok
\SourceText{Definition of Delta and Lemma lem:gamma-independent.}
\uses{mod:basic-charactertypes,mod:basic-characterconductors,mod:localfield-finitequotients,mod:basic-phase}
\FormalizationContract The integration domain is definitionally the
multiplicative group $U_F^0$, and the public local constant requires an
explicit positive left Haar measure and an admissible denominator.  The
measure remains unnormalised.  Measure independence is proved from a
strictly positive Haar-uniqueness scalar, while denominator independence is
proved only under the raw integral's explicit nonvanishing hypothesis.
The admissibility subtype is reducibly the same order-equality subtype used
by the later finite API, so their denominators and integrands agree
definitionally.
No finite Gauss sum or integral--finite comparison theorem is imported or
used.
\end{definition}
